\section{Requirements and Design Principles}

The requirements for the platform came from two directions. On the technology side, we needed
scalable ingestion, long retention, Kubernetes-native service discovery, log aggregation, alerting,
and dashboards. On the operations side, the need was harder to pin down but just as important: a
common language between teams. A graph is only useful when the right people understand what it
means and what to do next. On top of that, our infrastructure spans many types of hardware from
different vendors, each with its own quirks, and our users have very different expectations of what
the platform should show them. That diversity made the dashboard design particularly difficult;
there was no single answer that worked for everyone.

The DevOps operations rely on a set of core design principles, and any observability stack we built
had to reflect them from the ground up \cite{lsst-it-k8s-cookbook,prometheus-operator,grafana-loki-alloy,grafana-loki}.

\begin{center}
\begin{minipage}{\linewidth}
\renewcommand{\arraystretch}{1.3}
\small
\textbf{Table 1. Core Design Principles and Their Operational Rationale}\\[4pt]
\begin{tabular}{p{0.22\linewidth} p{0.36\linewidth} p{0.32\linewidth}}
\hline
\textbf{Principle} & \textbf{Implementation pattern in the platform} & \textbf{Why it matters for operations} \\
\hline
\textbf{Configuration as code} &
  Keep dashboards, alert rules, Helm values, and overlays in Git. &
  Changes are reviewed, promoted, and recovered like the rest of the infrastructure. \\[4pt]
\textbf{Open-source first} &
  Use Prometheus, Grafana, Loki, Mimir, Alloy, Fleet, exporters, and Kubernetes ecosystem tools. &
  The team can inspect the system, tune it, and avoid a single vendor dependency. \\[4pt]
\textbf{Kubernetes-native discovery} &
  Use ServiceMonitor, PodMonitor, Probe, and label selectors. &
  Teams can onboard services using conventions rather than one-off manual edits. \\[4pt]
\textbf{Multi-site and multi-cluster awareness} &
  Carry labels such as \texttt{site} and \texttt{cluster} into metrics and logs. &
  Operators can compare behavior across environments and avoid guessing where a signal came from. \\[4pt]
\textbf{Logs and metrics together} &
  Collect Kubernetes logs, node logs, events, and metrics into systems that can be queried from Grafana. &
  Incidents usually need both time-series trends and exact log messages. \\[4pt]
\textbf{Room for constant improvement} &
  Treat dashboards and alerts as living assets, not as a finished project. &
  Commissioning and operations create new questions, so the platform must keep changing. \\
\hline
\end{tabular}
\end{minipage}
\end{center}